% Generated draft: requires mathematical and source review.
\documentclass[11pt]{article}
\usepackage[margin=25mm]{geometry}
\usepackage{fontspec}
\setmainfont{Latin Modern Roman}
\usepackage{amsmath,amssymb,mathtools,unicode-math}
\setmathfont{Latin Modern Math}
\usepackage{enumitem,needspace,xurl,hyperref}
\setlength{\parindent}{0pt}
\setlength{\parskip}{5pt}
\setlength{\emergencystretch}{3em}
\newcommand{\N}{\mathbb{N}}
\newcommand{\Z}{\mathbb{Z}}
\newcommand{\Q}{\mathbb{Q}}
\newcommand{\R}{\mathbb{R}}
\newcommand{\C}{\mathbb{C}}
\newcommand{\F}{\mathbb{F}}
\newcommand{\A}{\mathbb{A}}
\newcommand{\PP}{\mathbb{P}}
\newcommand{\E}{\mathbb{E}}
\newcommand{\eps}{\varepsilon}
\newcommand{\dd}{\,\mathrm{d}}
\begin{document}
\section*{UnsolvedMath ID: 800188}
\textbf{Source problem code:} catalogue entry
\subsection{Definitions and mathematical statement}
For an integer \(k\geq 3\) and a positive integer \(m\), define the \(k\)-fold divisor function by
\[
d_k(m)=\#\{(a_1,\ldots,a_k)\in\Z_{\geq1}^{k}:a_1\cdots a_k=m\},
\qquad d(m)=d_2(m).
\]
For real \(T>1\) and \(\epsilon_0\geq0\), put
\[
I_{k,\epsilon_0}(T)=\int_0^T
\left(\sum_{\substack{m\in\Z_{\geq1}\\m\leq T^{2+\epsilon_0}}}\frac{d_k(m)}{m^{1/2+it}}\right)
\left(\sum_{\substack{n\in\Z_{\geq1}\\n\leq T^{2+\epsilon_0}}}\frac{d(n)}{n^{1/2-it}}\right)\,\dd t,
\]
where \(i^2=-1\) and \(x^z=\exp(z\log x)\) for positive real \(x\). The sums therefore end at \(\lfloor T^{2+\epsilon_0}\rfloor\).

\textbf{Problem.} For each fixed integer \(k\geq3\), asymptotically evaluate \(I_{k,\epsilon_0}(T)\) as \(T\to\infty\), for some small \(\epsilon_0\geq0\). The parameter \(\epsilon_0\) is understood to be fixed as \(T\) varies. The source does not specify the meaning of ``small,'' whether its choice may depend on \(k\), a proposed main term, or a required error bound. Its stated inequality permits \(\epsilon_0=0\); it does not explicitly require a positive value. The displayed integral is the target as written, without taking a real part or an absolute value.

The source additionally states that mean values of long Dirichlet polynomials of length \(N=T^{2-\epsilon}\) can be evaluated conjecturally on the additive divisor conjecture, citing [HN]. This is supplied conditional background, not an independently verified theorem here. Neither the precise additive divisor conjecture nor the identity of [HN] is included in the record, and the problem does not specify whether its requested evaluation may assume that conjecture.
\subsection{Short English statement}
For each fixed integer \(k\geq3\), find an asymptotic formula for the integral of the two displayed divisor Dirichlet polynomials, both truncated at \(T^{2+\epsilon_0}\), as \(T\to\infty\), for some small fixed \(\epsilon_0\geq0\). The admissible parameter range, error requirement, and permitted conjectural assumptions need clarification.
\subsection{Sources}
\begin{itemize}
\item Local UnsolvedMath registry, permanent ID \(800188\), source problem code \texttt{catalogue entry}, title ``Moments of L-functions --- Moments of the Riemann zeta function --- Problem 40.'' The supplied registry labels the problem open; this status has not been independently verified.
\item Hamieh, Kadiri, Ng, \emph{Moments of L-functions Problem List} (2024), Problem 40, as attributed in the supplied record: \url{https://arxiv.org/abs/2405.17800}. The supplied provenance records access on 2026-10-16 and extraction by \texttt{arxiv\_latex}; these are supplied provenance details, not a claim of independent source access or verification.
\item [HN], cited in the supplied statement for the conditional evaluation at length \(N=T^{2-\epsilon}\). No bibliographic identification or separate URL for [HN] is supplied.
\end{itemize}
\end{document}
